\documentclass[11pt,a4paper,sans]{moderncv}
\moderncvstyle{classic}
\moderncvcolor{blue}
\usepackage[T1]{fontenc}
\usepackage{lmodern}
\usepackage[utf8]{inputenc}
\usepackage[polish]{babel}
\usepackage[scale=0.82]{geometry}
\name{Sherlock}{Holmes}
\title{Kandydat na analityka kryminalistycznego}
\address{221B Baker Street}{Londyn}{}
\phone[mobile]{+44 7700 900 123}
\email{sherlock.holmes@example.org}
\extrainfo{CV fikcyjne, przygotowane do ćwiczenia z LaTeX-a}
\begin{document}
\makecvtitle
\section{Profil}
\cvitem{}{Analityk specjalizujący się w rekonstrukcji zdarzeń na podstawie
śladów materialnych i zeznań. Łączy obserwację terenową z weryfikacją
hipotez; wyniki przedstawia w krótkich, sprawdzalnych raportach.}
\section{Doświadczenie}
\cventry{1881--1904}{Konsultant dochodzeniowy}{Praktyka własna}
{Londyn}{}{Analiza miejsc zdarzeń, porównywanie śladów, budowa osi czasu
i przedstawianie wniosków klientom oraz organom śledczym.}
\cventry{1881--1904}{Współpraca z lekarzem}{Praktyka własna}
{Londyn}{}{Konsultowanie hipotez medycznych i dokumentowanie obserwacji
w formie nadającej się do niezależnej kontroli.}
\section{Wybrane kompetencje}
\cvitem{Analiza}{Dedukcja, ocena wiarygodności źródeł, testowanie
alternatywnych wyjaśnień.}
\cvitem{Laboratorium}{Podstawowe badania chemiczne i mikroskopowa
obserwacja materiału dowodowego.}
\cvitem{Komunikacja}{Raporty z rozróżnieniem obserwacji, hipotez i wniosków.}
\section{Przykładowy projekt}
\cvitem{Rekonstrukcja}{Ustalenie kolejności zdarzeń z niezależnych
zeznań i śladów przy zachowaniu jawnych założeń oraz niepewności.}
\section{Języki}
\cvitemwithcomment{Angielski}{ojczysty}{}
\cvitemwithcomment{Francuski}{roboczy}{}
\section{Informacja o dokumencie}
\cvitem{}{To artystyczna interpretacja postaci Sherlocka Holmesa,
bohatera prozy Arthura Conana Doyle'a. Dane kontaktowe i opis stanowiska
są fikcyjne. Skład oparto na klasie \texttt{moderncv} (styl \texttt{classic}).}
\end{document}
